% Generated draft: requires mathematical and source review.
\documentclass[11pt]{article}
\usepackage[margin=25mm]{geometry}
\usepackage{fontspec}
\setmainfont{Latin Modern Roman}
\usepackage{amsmath,amssymb,mathtools,unicode-math}
\setmathfont{Latin Modern Math}
\usepackage{enumitem,needspace,xurl,hyperref}
\setlength{\parindent}{0pt}
\setlength{\parskip}{5pt}
\setlength{\emergencystretch}{3em}
\newcommand{\N}{\mathbb{N}}
\newcommand{\Z}{\mathbb{Z}}
\newcommand{\Q}{\mathbb{Q}}
\newcommand{\R}{\mathbb{R}}
\newcommand{\C}{\mathbb{C}}
\newcommand{\F}{\mathbb{F}}
\newcommand{\A}{\mathbb{A}}
\newcommand{\PP}{\mathbb{P}}
\newcommand{\E}{\mathbb{E}}
\newcommand{\eps}{\varepsilon}
\newcommand{\dd}{\,\mathrm{d}}
\begin{document}
\section*{UnsolvedMath ID: 20003170}
\textbf{Source problem code:} catalogue entry
\subsection{Definitions and mathematical statement}
Let $G$ be a compact Lie group with a left-invariant Riemannian metric and $H<G$ a closed subgroup. Left multiplication by $H$ is isometric and free, so the left quotient $H\backslash G$ has its Riemannian-submersion metric. Quasi-positive curvature means nonnegative sectional curvature everywhere and strictly positive curvature on every two-plane at at least one point. It is different from positivity on an open dense set. The source asks for new examples of such quotients but does not define new, restrict $G,H$, or state a nonnegative-curvature hypothesis on the metric upstairs.

The aid's no-new result fixes connected $G$, a closed (not necessarily connected) subgroup $H$, and
\[
 [\mathfrak g,\mathfrak g]\cong\mathfrak{su}(2),\qquad
 \dim(H\backslash G)\geq2.
\]
Write $\mathfrak z$ for the center of $\mathfrak g$ and $\operatorname{pr}_{\mathfrak z}$ for projection along the derived ideal. It reports equivalence of
\[
 \begin{gathered}
 H\backslash G\text{ admits some quasi-positively curved metric},\\
 \pi_1(H\backslash G)\text{ is finite},\\
 \operatorname{pr}_{\mathfrak z}(\mathfrak h)=\mathfrak z.
 \end{gathered}
\]
The first condition concerns any metric, hence includes every possible quasi-positive left-invariant-quotient construction on this underlying space.

When these conditions hold, the stronger reported conclusion is that \emph{every} bi-invariant metric on $G$ induces strictly positive normal homogeneous sectional curvature on the same quotient. In the standard normal-homogeneous convention, if $\mathfrak m=\mathfrak h^\perp$, the sectional numerator for $X,Y\in\mathfrak m$ is
\[
 \|[X,Y]_{\mathfrak h}\|^2+
 \frac14\|[X,Y]_{\mathfrak m}\|^2.
\]
The relevant orthogonality and bracket normalizations must match the chosen bi-invariant metric. The source gives no proof or full reference for the equivalence and its application to this formula.

Consequently the aid excludes new underlying quasi-positive Geroch examples in this rank-one-derived class, including $U(2)$ and the corresponding torus-by-$SU(2)$ groups: whenever such a quotient admits quasi-positive curvature, it already has positive normal homogeneous curvature. This does not prohibit new metric formulas on an already known space. ``Rank-one-derived'' refers to the derived Lie algebra and does not require the full group to have rank one.

The summary additionally reports a persistent-central-plane algebraic obstruction for arbitrary compact $G$, but does not state its criterion or hypotheses. No explicit obstruction formula is invented from its name, and the rank-one theorem is not extended to arbitrary higher derived rank. The higher-rank existence question remains open in the supplied aid; the original general search for new quotients is not settled by the special class.
\subsection{Short English statement}
The reported theorem says that in the specified $SU(2)$-derived class, every quotient supporting any quasi-positive metric already has positive normal homogeneous metrics. This rules out new underlying examples in that class, while leaving higher-rank groups and the unstated central-plane criterion outside the explicit result.
\subsection{Sources}
\begin{itemize}
\item Local UnsolvedMath registry, permanent ID 20003170, source code catalogue entry; \texttt{aim-topology-notes.json}, notes[257].
\item AIM, \emph{Manifolds with nonnegative sectional curvature}, Problem 45: \url{https://aimath.org/WWN/nnsectcurvature/nnsectcurvature.pdf}. No proposer, remarks or original bibliography is supplied.
\item Geroch quotient terminology, normal homogeneous curvature and the persistent-central-plane obstruction are invoked by the supplied summary without full bibliographic references or exact obstruction statement.
\item Machine-generated research attempt 1: partial result, open at the research run, rigor labelled ``proved,'' low novelty confidence; registry partially solved. Not peer reviewed; independently verify before citation. Registry updated 14 August 2026.
\end{itemize}
\end{document}
